\subsection*{Cancellation of the A.11 local-normalization derivatives}
\label{subsec:A11-local-derivative-cancellation}
On and after the two source (A.11) corrections, write
$E_c(y,\eta)=E_b(y)[1+u(y,\eta)]$, where
$u=\sum_{i=1}^2c_i(\eta)\beta(y-y_i)$ and
$E_b$ is independent of $\eta$. The two disjoint corrected supports
have total logarithmic length $3/5$. This permits an exact
cancellation in the axial source (4.16), stronger than the previous
separate estimates for the parameter derivatives of the tails.

\begin{theorem}[Exact future-support formula]\label{thm:A11-local-cancel}
Assume the \emph{genuine} source normalization
$S(\infty,\eta)=\Pi(\infty,\eta)=0$, and let
$D_c=S/(XE_c^2)$, $P_c=-\Pi/E_c^2$,
$g=\partial_\eta\log E_c$,
$R_b(y,t)=\exp[2\int_0^t\ell_b(y+v)\,dv]$.
Then on either A.11 support,
\begin{equation}\label{eq:A11-local-cancel}
2g(D_c-P_c)+\partial_\eta D_c-\partial_\eta P_c
=\frac{\int_0^\infty R_b(y,t)(1-e^{-t})
 [1+u(y+t)]u_\eta(y+t)\,dt}{[1+u(y)]^2}.
\end{equation}
The right-hand integrand vanishes outside the future portions
of the two corrected supports. If $\ell_b\le0$ at all future
radii, $0<\lambda\le1/20$,
$|c_i|\le192\lambda^{28}$,
$|c_{i,\eta}|\le12\lambda^{28}$ and
$\|\beta\|_\infty\le640/3$, then, with
$\delta=40960\lambda^{28}$,
\begin{equation}\label{eq:A11-local-cancel-bound}
|2g(D_c-P_c)+D_{c,\eta}-P_{c,\eta}|
\le1536\lambda^{28}\frac{1+\delta}{(1-\delta)^2}
<\frac1{1000}.
\end{equation}
\end{theorem}
\begin{proof}
Both source tails are positive integrals. At the current radial
position put $v=1+u(y)$ and write their unnormalized numerators
$A_D,A_P$, so $D_c=A_D/(2v^2)$ and
$P_c=A_P/(2v^2)$. Since the base field is $\eta$-independent,
$g=v_\eta/v$. Direct differentiation gives
$\partial_\eta(A/(2v^2))+2g A/(2v^2)=
A_\eta/(2v^2)$, canceling all local-normalization derivatives.
After changing radial integration variable $x=Xe^t$, the
difference $A_D-A_P$ has the kernel $R_b(y,t)(1-e^{-t})$
times $[1+u(y+t)]^2$, proving \eqref{eq:A11-local-cancel}.
Because $\ell_b\le0$, $0\le R_b\le1$; because the supports
are disjoint, $|u_\eta|\le2560\lambda^{28}$ and
their total future length is at most $3/5$. Divide by
$v^2\ge(1-\delta)^2$ to obtain the bound.
\end{proof}

\begin{corollary}[Improved strict stress and cone margins]
\label{cor:A11-local-cancel-cone}
Assume also $D_c<233/100$, $0<h\le\lambda/100$,
$E_c\le\lambda/4$ on both supports, and the future \emph{baseline}
slope is nonpositive. The source axial stress obeys
\begin{equation}\label{eq:A11-new-margin}
\frac{|N_s|}{E_c}<\frac{63}{250}\lambda,\qquad
\frac{\lambda}{2}-\frac{|N_s|}{E_c}>
\frac{31}{125}\lambda.
\end{equation}
This conclusion does not require a separate assumption that the
\emph{corrected} future slope be nonpositive. If additionally
$Q_s\ge\lambda/2$, $-3\lambda/2\le\ell_c\le-\lambda/2$
and $X\ge9/\lambda$, the normalized cone inequality satisfies
\begin{equation}\label{eq:A11-new-cone}
2(1-a/p_1)^2-(a-2)w^2
\ge\frac{144343}{312500}>\frac9{20},
\quad a=2-2\ell_c,\quad
p_1=\frac{XQ_s}{1-2h\eta^2},\quad
w=\frac{N_s}{E_cQ_s}.
\end{equation}
\end{corollary}
\begin{proof}
The future baseline slope gives $P_b\le1/2$; positive tail
comparison gives $P_c\le(1+8\delta)/2$. Thus
$4AP_c<1002/1000$, $4hD_c<1/200$ and the
derivative term is $<1/1000$. The source (4.16) formula
bounds $|N_s|/E_c$ by $E_c$ times their sum,
which is less than $E_c(1008/1000)\le(63/250)\lambda$.
The exact slack is $31\lambda/125$. With the
extra cone hypotheses, $a/p_1\le1/2$,
$a-2\le3/20$ and $|w|\le63/125$. Hence
$2(1-a/p_1)^2-(a-2)w^2
\ge1/2-(3/20)(63/125)^2=144343/312500$.
\end{proof}

The statements are **conditional mathematical implications**, not
a source-authenticated interval for one completely frozen Appendix-A
parameter tuple. The terminal moment/pressure gauges, coefficient
enclosures, complete exterior and pulse stress gates, nonlinear
Appendix-B profile, distinct unforced $n=1,2$ coefficient solutions
and the zero-force physical PDE remain unverified. The originating
OpenAI global Theorem~1.1 includes smooth forcing.